
\section{Lorentz Transforms}
Metric Tensors are used to define distance metrics in 4D spacetime, called spacetime intervals. $g_{\mu\nu}$ is the general metric tensor, and $\eta_{\mu\nu}$ is used for explicitly flat spacetime. For particle physics, either can be used, and it would be understood that $g_{\mu\nu}$ is in flat spacetime. 
\begin{align}
	g_{\mu\nu}&=
	\begin{bmatrix}
		g_{00} & g_{01} & g_{02} & g_{03} \\
		g_{10} & g_{11} & g_{12} & g_{13} \\
		g_{20} & g_{21} & g_{22} & g_{23} \\
		g_{30} & g_{31} & g_{32} & g_{33} \\
	\end{bmatrix}
	\\
	\eta_{\mu\nu}&=\mqty[\dmat[0]{1,-1,-1,-1}]
\end{align}
Four Vectors:
\begin{align*}
	x^{\mu}&=\colvec{ct\\x_{1}\\x_{2}\\x_{3}}=\colvec{ct\\\vb{x}}
	&
	p^{\mu}&=\colvec{cE\\p_{1}\\p_{2}\\p_{3}}=\colvec{cE\\\vb{p}}
	\\[0.5cm]
	A^{\mu}&=\colvec{\oo{c}\phi\\A_{1}\\A_{2}\\A_{3}}=\colvec{\oo{c}\phi\\\vb{A}}
	&
	\pder*[\mu]{}&=\colvec{\oo{c}\pder*{t}\\\pder*{1}\\\pder*{2}\\\pder*{3}}=\colvec{\oo{c}\pder*{t}\\\grad{}}
\end{align*}

Lorentz scalars are invariant under lorentz transformations, this is centrally important, and a Lorentz invariant Lagrangian/Action should be built on Lorentz scalars. This property is encapsulated in a 
frequently used identity, which is obtained by considering the inner product between two four vectors, which is a Lorentz scalar:
\begin{align*}
	a\cdot b&=a^{\mu}\eta_{\mu\nu}b^{\nu}
	\\
	\Big\downarrow\Lambda& \notag \\
	a^{\prime}\cdot b^{\prime}&=a^{\rho}\Lambda^{\mu}_{~\rho}\eta_{\mu\nu} \Lambda^{\nu}_{~\sigma}b^{\sigma}=a\cdot b=a^{\rho}\eta_{\rho\sigma}b^{\sigma}
	\\
	\Lambda^{\mu}_{~\rho}\eta_{\mu\nu} \Lambda^{\nu}_{~\sigma}a^{\rho}b^{\sigma}&=\eta_{\sigma\rho}a^{\rho}b^{\sigma}
\end{align*}
\begin{align}		
	\boxed{\Lambda^{\mu}_{~\rho}\eta_{\mu\nu} \Lambda^{\nu}_{~\sigma}=\eta_{\sigma\rho}}\label{eqn:LorentzMatrixConstraint}
\end{align}
This is a \emph{constraint} on the Lorentz transform matrix.
\\
\(D[\Lambda]^{a}_{~b\,}\) is the matrix representation of the Lorentz Group, and has the properties:
\begin{align*}
	D[\Lambda_1]D[\Lambda_2]&=D[\Lambda_1\Lambda_2]
	\\
	D[\Lambda^{-1}]&=D[\Lambda]^{-1}
	\\
	D[1]&=1
\end{align*}
To find the representations of a group, the typical formulation is to take infinitesimal transformations of the group and to study the resulting algebra. 
\begin{align*}
	\Lambda^{\mu}_{~\nu}=\underbrace{\delta^{\mu}_{~\nu}}_{\rm Identity}+\underbrace{\omega^{\mu}_{~\nu}}_{\rm Infinitesmial}
\end{align*}
From the constraint (\ref{eqn:LorentzMatrixConstraint}), it can be shown that \(\omega^{\mu}_{~\nu}\) must be anti-symmetric, which means there are 6 unique off-diagonal elements, which correspond to the 
\subsection{Definitions used in the field (of study):}
Lorentz transforms and causality
\\
spacetime intervals

\subsection{Lorentz Transformations of general field}
\begin{align}
	\phi^{a}(\vb{x})&\xrightarrow{\Lambda}\qty(\phi^{a}(\vb{x}))^{\prime}=D[\Lambda]^{a}_{~b\,}\phi^{b}(\Lambda^{-1}\vb{x})
\end{align}
Where \(D[\Lambda]^{a}_{~b\,}\) is the matrix representation of the Lorentz Group.
